% ═══════════════════════════════════════════════════════════════════════════════
%  CHAPITRE 1 : PRÉSENTATION DE L'ORGANISME D'ACCUEIL
% ═══════════════════════════════════════════════════════════════════════════════
\newpage
\section{Présentation de l'organisme d'accueil}
\label{ch1}

\subsection{Introduction}
\label{ch1:intro}

Ce chapitre présente l'environnement dans lequel s'est déroulé notre projet de fin d'études. Il décrit la Faculté des Sciences Ben M'Sick (FSBM) -- son histoire, ses missions, sa structure organisationnelle -- ainsi que le département d'accueil et le cadre spécifique du stage. Cette contextualisation permet de mieux comprendre les enjeux institutionnels et pédagogiques qui ont motivé la réalisation de notre système de gestion des soutenances.

\subsection{Présentation générale de l'organisme}
\label{ch1:presentation}

\subsubsection{Historique et fiche signalétique}
\label{ch1:historique}

La Faculté des Sciences Ben M'Sick (FSBM) est l'une des composantes de l'Université Hassan II de Casablanca, l'une des plus grandes universités du Maroc. Située dans l'arrondissement de Ben M'Sick, à Casablanca, elle a été créée pour répondre à la demande croissante de formation scientifique dans la région du Grand Casablanca.

\begin{table}[H]
    \centering
    \small
    \begin{tabular}{lp{10cm}}
        \toprule
        \rowcolor{primary!10}
        \textbf{Information} & \textbf{Détail} \\
        \midrule
        Nom officiel         & Faculté des Sciences Ben M'Sick (FSBM) \\
        Université de rattachement & Université Hassan II de Casablanca \\
        Adresse              & Avenue Commandant Driss El Harti, Ben M'Sick, Casablanca \\
        Année de création    & 1997 \\
        Type                 & Établissement public d'enseignement supérieur \\
        Domaines             & Sciences exactes, Sciences et technologies \\
        \bottomrule
    \end{tabular}
    \caption{Fiche signalétique de la FSBM}
    \label{tab:fsbm_fiche}
\end{table}

\subsubsection{Missions et activités}
\label{ch1:missions}

La FSBM remplit plusieurs missions fondamentales. Elle assure tout d'abord la formation initiale en dispensant un enseignement supérieur de qualité dans les disciplines scientifiques fondamentales et appliquées aux niveaux Licence, Master et Doctorat. Parallèlement, elle développe une activité intense de recherche scientifique au sein de ses divers laboratoires et équipes. L'établissement se consacre également à préparer activement ses étudiants à leur insertion dans le monde professionnel à travers des formations adaptées, tout en contribuant au rayonnement académique, scientifique et technologique de la région et du pays.

\subsubsection{Secteurs d'intervention}
\label{ch1:secteurs}

La faculté couvre un spectre étendu de disciplines scientifiques, incluant les mathématiques—avec l'algèbre, l'analyse, les probabilités, les statistiques et les mathématiques appliquées—ainsi que l'informatique, qui englobe le développement logiciel, les bases de données, les réseaux et l'intelligence artificielle. Les sciences physiques, la chimie, les sciences de la vie (biologie moléculaire, biotechnologie, écologie) et les sciences de la terre (géologie, géophysique, géomatique) complètent cette offre de formation diversifiée.

\subsection{Structure organisationnelle}
\label{ch1:organisation}

\subsubsection{Organigramme}
\label{ch1:organigramme}

La FSBM est administrée par un Doyen, assisté par des Vice-Doyens et un Secrétaire Général. Sa structure s'articule autour de plusieurs organes décisionnels et entités fonctionnelles, au premier rang desquels figure le Conseil de faculté, instance délibérante composée d'enseignants-chercheurs, d'étudiants et de représentants administratifs. Le fonctionnement académique est encadré par des commissions pédagogiques chargées de l'organisation et du suivi des filières. La faculté s'appuie également sur des départements pédagogiques, des laboratoires de recherche ainsi que des services administratifs essentiels, notamment la scolarité, les ressources humaines, les finances et la bibliothèque.

\subsubsection{Description des départements}
\label{ch1:departements}

La structuration pédagogique repose sur plusieurs départements, chacun responsable d'une ou plusieurs filières. Le Département de Mathématiques propose les licences en Mathématiques Fondamentales et Appliquées, en plus des masters spécialisés. Le Département d'Informatique est responsable de la filière Licence MIP (Mathématiques, Informatique et Physique) et de plusieurs masters en informatique. Les autres domaines sont couverts par le Département de Physique, le Département de Chimie, le Département de Biologie et le Département de Géologie, qui offrent chacun des formations spécialisées. Ces départements sont dirigés par des chefs élus, assistés par un conseil de département, assurant la coordination des enseignements et des emplois du temps.

\subsection{Environnement du stage}
\label{ch1:environnement}

\subsubsection{Département d'accueil}
\label{ch1:equipe}

Notre projet de fin d'études s'est déroulé au sein de la filière \textbf{Licence Fondamentale MIP (Mathématiques, Informatique et Physique)} du Département d'Informatique. Cette filière vise à former des étudiants polyvalents possédant une solide base en mathématiques et une maîtrise des technologies de l'information. La formation couvre des domaines variés tels que le développement logiciel, les bases de données, les réseaux, l'algorithmique avancée et l'intelligence artificielle.

Le corps enseignant de la filière est composé de professeurs de l'enseignement supérieur, de professeurs habilités et de professeurs assistants, couvrant l'ensemble des spécialités du programme.

\subsubsection{Contexte du projet au sein de l'organisme}
\label{ch1:contexte_projet}

Le projet de système de gestion des soutenances et des jurys s'inscrit dans une volonté de modernisation des processus administratifs et pédagogiques au sein de la faculté. Actuellement, l'organisation des soutenances -- qu'il s'agisse des soutenances de PFE, de mémoires de master ou de projets de licence -- repose sur des méthodes manuelles et des outils bureautiques génériques.

Le département MIP, confronté à un volume croissant d'étudiants et à une complexité accrue dans la coordination des jurys, a identifié le besoin d'un outil informatique dédié. Ce projet a donc été proposé par \textbf{Pr. Mohammed AIT DAOUD} dans le cadre des projets de fin d'études de la licence MIP, avec pour objectif de fournir une solution concrète et opérationnelle répondant aux besoins réels du département.

\subsection{Conclusion}
\label{ch1:conclusion}

La Faculté des Sciences Ben M'Sick constitue un cadre académique riche et diversifié, offrant des formations scientifiques de qualité dans un large éventail de disciplines. Le département MIP, par sa proximité avec les technologies de l'information, est naturellement le lieu idéal pour le développement d'une solution numérique de gestion des soutenances. Ce contexte a directement inspiré les choix technologiques et architecturaux qui seront présentés dans les chapitres suivants.
